\begin{tikzpicture}[figure picture,x=1mm,y=1mm,
  font=\figurefont\fontsize{8}{10}\selectfont,
  state/.style={circle,fill=FigureInk,inner sep=1.1pt},
  link/.style={draw=FigureStroke,line width=.65pt}]
  \path[use as bounding box] (0,0) rectangle (165,69);
  \node at (26,65) {第二步已忘记第一步};
  \node[state] (root) at (26,51) {};
  \node[above=1mm of root] {\(I_1\)};
  \node[state] (left) at (12,33) {};
  \node[state] (right) at (40,33) {};
  \draw[link] (root)--node[pos=.5,left] {\(L\)} (left);
  \draw[link] (root)--node[pos=.5,right] {\(R\)} (right);
  \draw[FigureMuted,dashed,rounded corners=1.8pt]
    (8,29.5) rectangle (44,36.5);
  \node[fill=white,inner sep=1pt] at (26,33) {\(I_2\)};
  \coordinate (Ll) at (4,11);
  \coordinate (Lr) at (20,11);
  \coordinate (Rl) at (32,11);
  \coordinate (Rr) at (48,11);
  \draw[link] (left)--node[pos=.5,left] {\(l\)} (Ll);
  \draw[link] (left)--node[pos=.5,right] {\(r\)} (Lr);
  \draw[link] (right)--node[pos=.5,left] {\(l\)} (Rl);
  \draw[link] (right)--node[pos=.5,right] {\(r\)} (Rr);
  \foreach \name in {Ll,Lr,Rl,Rr} {
    \fill[FigureInk] (\name) circle (.55mm);
    \node[below=1mm] at (\name) {\(\name\)};
  }
  \node at (84,65) {开局混合完整计划};
  \node at (137,65) {逐处独立随机化};
  \node at (84,57) {\(\mu(L,l)=\mu(R,r)=1/2\)};
  \node at (137,57) {\(p=q=1/2\)};
  \foreach \shift in {61,114} {
    \begin{scope}[xshift=\shift mm]
      \node at (16.5,47) {\(l\)};
      \node at (29.5,47) {\(r\)};
      \node at (5,37.5) {\(L\)};
      \node at (5,24.5) {\(R\)};
      \fill[white] (10,18) rectangle (36,44);
      \ifnum\shift=61
        \fill[FigureBlueFill] (10,31) rectangle (23,44);
        \fill[FigureBlueFill] (23,18) rectangle (36,31);
        \node at (16.5,37.5) {\(0.5\)};
        \node at (29.5,37.5) {\(0\)};
        \node at (16.5,24.5) {\(0\)};
        \node at (29.5,24.5) {\(0.5\)};
      \else
        \fill[FigureGreenFill] (10,18) rectangle (36,44);
        \foreach \xx in {16.5,29.5}
          \foreach \yy in {24.5,37.5}
            \node at (\xx,\yy) {\(0.25\)};
      \fi
      \draw[FigureGrid] (10,18) rectangle (36,44);
      \draw[FigureGrid] (23,18)--(23,44);
      \draw[FigureGrid] (10,31)--(36,31);
    \end{scope}
  }
  \node[font=\FigureNoteFont\fontsize{7.5}{9}\selectfont] at (84,8)
    {随机计划保留跨时相关};
  \node[font=\FigureNoteFont\fontsize{7.5}{9}\selectfont] at (137,8)
    {还会产生 \(Lr\) 与 \(Rl\)};
\end{tikzpicture}
